\documentclass[11pt,a4paper]{article}
\usepackage{xeCJK}
\usepackage{fontspec}
\usepackage{xcolor}
\usepackage{fancyhdr}
\usepackage{booktabs}
\usepackage{array}
\usepackage{amssymb}
\usepackage{geometry}
\usepackage[protrusion=true]{microtype}
\geometry{margin=1.8cm,top=2cm,headheight=15pt}
\linespread{1.05}

\setmainfont{Baskerville}
\setCJKmainfont{Songti SC}
\newfontfamily{\starfont}{Songti SC}

\definecolor{wrongred}{RGB}{198,40,40}
\definecolor{rulegray}{RGB}{200,200,200}

\setlength{\emergencystretch}{3em}
\setlength{\fboxsep}{3pt}
\setlength{\parindent}{0pt}

\newcommand{\checkmarkbox}{\fbox{\rule{0pt}{0.8ex}\rule{0.8ex}{0pt}}}
\newcommand{\STAR}{\textcolor{wrongred}{\starfont ★}}
\newcommand{\bk}{\rule[-0.55ex]{1.0em}{0.4pt}}
\newcommand{\A}[1]{\textbf{#1}}
\newcolumntype{L}[1]{>{\raggedright\arraybackslash}p{#1}}

\pagestyle{fancy}
\fancyhf{}
\fancyhead[L]{\footnotesize 英语翻译专项 V. Translation（19\% · 订正）}
\fancyhead[R]{\footnotesize 赵文瀚}
\fancyfoot[C]{\thepage}
\renewcommand{\headrulewidth}{0.4pt}
\renewcommand{\footrulewidth}{0pt}

\begin{document}

\begin{center}
{\LARGE\bfseries 英语翻译专项 · 订正}\\[0.25em]
{\large V. Translation（共 8 题 · 19 分）· 考点复盘}\\[0.35em]
{\footnotesize 2029 届高一英语 · 翻译（中 $\to$ 英）· 原卷 8 题（67--74）· 2026-10}
\end{center}

\vspace{-0.2em}
\noindent{\color{rulegray}\rule{\textwidth}{0.8pt}}

\vspace{0.5em}
\noindent\fcolorbox{black!15}{black!4}{%
  \parbox{\dimexpr\linewidth-2\fboxsep-2\fboxrule\relax}{%
    \small 日期：2026-10-11\qquad 来源：原卷照片 V. Translation（8 题 67--74，共 19 分）\qquad 姓名：赵文瀚\\[0.3em]
    题量与分值：\textbf{1.5$'$}$\times$4（67--70）\enspace+\enspace \textbf{3$'$}$\times$2（71--72）\enspace+\enspace \textbf{3.5$'$}$\times$2（73--74）\enspace=\enspace 19 分\qquad
    得分：\underline{\hspace{1.8cm}}\,/\,19\qquad
    重做日期：\underline{\hspace{2.4cm}}}}

\vspace{0.6em}
\noindent{\footnotesize 说明：原卷\textbf{仅印题干}，\textbf{无学生作答、无老师红笔}（页眉手写「642」，疑为作业号 260642，与 grammar\_03 同源；未印校名与卷名），因此本表\textbf{不设「我的作答／判定」栏}。下表按原卷题号 67--74 全列，参考答案按标准语法补拟，\A{加粗}＝与括号内提示词对应的答案部分。取得批改件后可另附一页逐题照录判定。}

\vspace{0.4em}
\renewcommand{\arraystretch}{1.15}
\noindent{\footnotesize\begin{tabular}{@{}L{1.35cm}@{\hspace{0.1cm}}L{4.85cm}@{\hspace{0.1cm}}L{6.4cm}@{\hspace{0.1cm}}L{4.35cm}@{}}
\toprule
\textbf{题号（分）} & \textbf{中文原句（提示词）} & \textbf{参考答案} & \textbf{考点 · 易错点} \\
\midrule
67 \newline {\scriptsize 1.5$'$} & 老师必须记录学生的出勤情况。\newline {\scriptsize（record n.）} & Teachers must \A{keep a record of} students' attendance. & record 作名词的搭配：keep a record of；attendance 不可数 \\
68 \newline {\scriptsize 1.5$'$} & 请允许我代表我的同事们感谢你主动提出帮忙。\newline {\scriptsize（behalf）} & Please allow me to thank you for your offer of help \A{on behalf of my colleagues}. & on behalf of＝代表（固定介词短语，勿写 in behalf of） \\
69 \newline {\scriptsize 1.5$'$} & 他定期向他的老板作业务汇报。\newline {\scriptsize（basis）} & He gives business reports to his boss \A{on a regular basis}. & on a regular basis＝定期（介词用 on，不用 in） \\
70 \newline {\scriptsize 1.5$'$} & 毫无疑问，人们的生活非常不同。\newline {\scriptsize（doubt）} & \A{There is no doubt that} people's lives are very different. & There is no doubt that＋陈述句；肯定句 doubt 后才用 whether \\
71 \newline {\scriptsize 3$'$} & 纸张是厚是薄无关紧要。\newline {\scriptsize（matter）} & \A{It doesn't matter whether} the paper is thick or thin. & it 作形式主语＋whether 从句（不用 if）；matter 不及物 \\
72 \newline {\scriptsize 3$'$} & 显然，自然灾害是造成该国经济危机的原因之一。\newline {\scriptsize（contribute）} & Obviously, natural disasters are one of the factors that \A{contribute to} the country's economic crisis. & contribute to＋名词；「原因之一」用 one of the factors（复数） \\
73 \newline {\scriptsize 3.5$'$} & 船上既拥挤又不舒适，同时疾病在船上传播。\newline {\scriptsize（spread）} & It was crowded and uncomfortable on the ship, and at the same time disease \A{spread} on board. & spread 过去式＝过去分词＝spread（勿写 spreaded）；on board＝在船上 \\
74 \newline {\scriptsize 3.5$'$} & 你没必要一早就动身，到那儿要不了多长时间。\newline {\scriptsize（take）} & \A{There is no need} for you to set off so early；it \A{won't take long} to get there. & there is no need (for sb) to do；take 表花费：it takes sb some time to do \\
\bottomrule
\end{tabular}}

\newpage
\begin{center}
{\Large\bfseries 考点归类与复盘}
\end{center}
\vspace{-0.4em}
\noindent{\color{rulegray}\rule{\textwidth}{0.8pt}}

\vspace{0.5em}
\noindent 说明：8 题按考点归类，供汇总对账；逐题参考答案见第 1 页。

\vspace{0.5em}
\renewcommand{\arraystretch}{1.3}
\begin{center}
\begin{tabular}{@{}p{9.6cm} p{4.2cm} c@{}}
\toprule
\textbf{考点} & \textbf{题号} & \textbf{题数} \\
\midrule
介词短语（on behalf of／on a regular basis） & {\small 68, 69} & 2 \\
名词搭配（keep a record of） & {\small 67} & 1 \\
固定句型（There is no doubt that…） & {\small 70} & 1 \\
形式主语 it＋whether 从句（matter） & {\small 71} & 1 \\
动词短语（contribute to＋one of the factors） & {\small 72} & 1 \\
词形（spread 过去式／过去分词同形） & {\small 73} & 1 \\
不定式句型（there is no need to do；take 表花费） & {\small 74} & 1 \\
\midrule
\textbf{合计} & \textbf{8 题} & \textbf{8} \\
\bottomrule
\end{tabular}
\end{center}

\vspace{0.6em}
\noindent\textbf{易错点提示}
\vspace{0.25em}

\noindent\begin{tabular}{@{}p{8.2cm}@{\hspace{0.4cm}}p{8.2cm}@{}}
\begin{minipage}[t]{\linewidth}
\small
\STAR\ \textbf{record}：名词用 \textbf{keep a record of}（勿写 record of 单独作谓语）；\emph{attendance} 不可数\\[0.2em]
\STAR\ \textbf{behalf}：只用于 \textbf{on behalf of}（代表），无 \emph{in behalf of}\\[0.2em]
\STAR\ \textbf{basis}：\textbf{on} a regular／daily \textbf{basis}（定期），介词是 on\\[0.2em]
\STAR\ \textbf{doubt}：\textbf{There is no doubt that}…（无疑）；否定句接 that，肯定句接 whether／if
\end{minipage} &
\begin{minipage}[t]{\linewidth}
\small
\STAR\ \textbf{matter}：\textbf{It doesn't matter whether}…（不用 if；whether 引导主语从句）\\[0.2em]
\STAR\ \textbf{contribute}：\textbf{contribute to} sth.；「……之一」用 \textbf{one of the factors}（复数）\\[0.2em]
\STAR\ \textbf{spread}：过去式＝过去分词＝\textbf{spread}（勿写 \emph{spreaded}）；\textbf{on board}＝在船上\\[0.2em]
\STAR\ \textbf{take}：\textbf{it takes sb some time to do}；\textbf{there is no need (for sb) to do}（没必要）
\end{minipage} \\
\end{tabular}

\vspace{0.65em}
\noindent\fcolorbox{black!15}{black!4}{%
  \parbox{\dimexpr\linewidth-2\fboxsep-2\fboxrule\relax}{%
    \textbf{复盘记录}\\[0.35em]
    最薄弱的板块：\underline{\hspace{9.0cm}}\\[0.55em]
    主要错误原因：\checkmarkbox 介词短语记错（on／in）\quad \checkmarkbox 固定句型不熟\quad \checkmarkbox 词形（spread／contribute）\quad \checkmarkbox 中文句式转换困难\\[0.55em]
    本次复习的改进措施：\underline{\hspace{10.1cm}}\\[0.55em]
    \hspace{3.2em}\underline{\hspace{10.1cm}}\\[0.55em]
    下次复习日期：\underline{\hspace{2.2cm}}\qquad
    当前掌握度：\checkmarkbox 仍薄弱\quad \checkmarkbox 基本掌握\quad \checkmarkbox 已掌握
  }}

\end{document}
